\subsection{Equivalence of DCP tableaux}\label{app:tableau-equivalence}

The stabilizing closure introduced in
\Cref{sec:dcp-stabilizer-formalism} identifies two DCP-stabilizer groups
precisely when they have the same fixed space. Here we give a direct
polynomial-time test for this equivalence in the pure-state case.

Let \(G_a\leq\mathcal D^n\) be a DCP-stabilizer group for
\(\ket{\psi_a}\), \(a\in\{1,2\}\), and retain a support representative
\(x_a\in\supp(\psi_a)\). Thus the data to be compared are
\begin{equation}
    T_a=(G_a,x_a),
    \qquad
    \Fix(G_a)=\C\ket{\psi_a}.
\end{equation}
Such a representative is already part of the phase-polynomial description
in \cref{eq:simulation-representation}; it may equally be carried
alongside the canonical group representation. There is no need for separate
update rules: its affine support is updated by the rules of
\Cref{sec:simulator-phase-polynomial-representation}.

Set \(V_a=X(G_a)\). Since the fixed space of \(G_a\) is one-dimensional,
\Cref{thm:atomic-basis} implies
\begin{equation}
    \supp(\psi_a)=x_a\oplus V_a.
    \label{eq:tableau-affine-support}
\end{equation}
Consequently the supports agree if and only if
\begin{equation}
    V_1=V_2
    \quad\text{and}\quad
    x_1\oplus x_2\in V_1.
    \label{eq:tableau-support-test}
\end{equation}

Suppose this test passes, write the common space as \(V\), and choose a basis
\(v^{(1)},\ldots,v^{(r)}\) of \(V\). For every \(i\), choose lifts
\begin{equation}
    g_{a,i}\in G_a,
    \qquad
    \pi(g_{a,i})=v^{(i)},
\end{equation}
and define the diagonal quotient
\begin{equation}
    \delta_i:=g_{1,i}g_{2,i}^{-1}\in\ker\pi.
\end{equation}
The lifts can be obtained by row reduction of the \(X\)-parts while carrying
the corresponding DCP operators.

\begin{proposition}[Tableau-equivalence criterion]
    \label{prop:tableau-equivalence}
    The tableaux \(T_1\) and \(T_2\) represent the same state up to global
    phase if and only if \cref{eq:tableau-support-test} holds and every
    $\delta_i$ acts as the identity on the common support
    $S=x_1\oplus V$.
\end{proposition}

\begin{proof}
    If the states agree up to global phase, their supports agree, giving
    \cref{eq:tableau-support-test}. Moreover, both \(g_{1,i}\) and
    \(g_{2,i}\) fix the same state. Their quotient is diagonal, so it must
    act as the identity on every computational-basis vector with nonzero
    amplitude, hence on all of \(S\).

    Conversely, suppose the stated conditions hold and write
    \begin{equation}
        \ket{\psi_a}=\sum_{x\in S}\alpha_a(x)\ket{x}.
    \end{equation}
    Decompose each lift as \(g_{a,i}=D_{a,i}X(v^{(i)})\), with \(D_{a,i}\)
    diagonal, and write
    \(D_{a,i}\ket{x}=d_{a,i}(x)\ket{x}\). The stabilizer equations give
    \begin{equation}
        \alpha_a(x)=d_{a,i}(x)\alpha_a(x\oplus v^{(i)}).
    \end{equation}
    Since
    \(\delta_i=D_{1,i}D_{2,i}^{-1}\) is the identity on \(S\), we have
    \(d_{1,i}(x)=d_{2,i}(x)\) there. Hence
    \(\alpha_1(x)/\alpha_2(x)\) is invariant under translation by every
    \(v^{(i)}\). These translations connect all points of \(S\), so this ratio
    is constant and \(\ket{\psi_1}\) and \(\ket{\psi_2}\) differ only by a
    global phase.
\end{proof}

It remains only to make the identity test explicit. Write
\begin{equation}
    \delta_i=\omega_8^{k_i}D(s_i,M_i),
\end{equation}
and let \(B\) be the matrix with columns \(v^{(1)},\ldots,v^{(r)}\). Its phase on a
basis vector is the Clifford phase polynomial
\begin{equation}
    F_i(x)
    =k_i+2s_i\cdot x
      +4\sum_{j<\ell}(M_i)_{j\ell}x_jx_\ell
      \pmod 8.
\end{equation}
By \Cref{cor:phase-polynomial-affine-pullback,thm:clifford-phase-polynomial},
the restriction \(c\mapsto F_i(x_1\oplus Bc)\) has Boolean degree at most
two, with its linear and quadratic coefficients divisible by \(2\) and \(4\),
respectively. It is therefore identically zero if and only if it vanishes at
\begin{equation}
    c=0,
    \qquad c=e_j,
    \qquad c=e_j\oplus e_\ell
    \quad (1\leq j<\ell\leq r).
    \label{eq:tableau-phase-test-points}
\end{equation}
Thus each diagonal quotient requires only \(O(r^2)\) phase evaluations. The
support comparison, construction of the lifts, and all phase tests are
polynomial-time operations. The result is independent of the chosen lifts:
two lifts of the same \(v^{(i)}\) differ by a diagonal stabilizer, which acts
identically on the corresponding support.
